% SOURCE_AUTHORITY: sga5_fr_workpass.tex, lines 10302--10314 (LF-normalized unit).
% SOURCE_SHA256: 791F4EFFC5E02832D5D77ED03518C8156D6F07E4C8238B03545DB93D883FBB28
\begin{statement}{Observação 6.3}
Como todo $S$-grupo redutivo é localmente, para a topologia étale, a imagem inversa pelo morfismo canônico $S\to\operatorname{Spec}(\mathbb Z)$ de um $\operatorname{Spec}(\mathbb Z)$-esquema redutivo (SGA 3 XXIII 5.7 e XXV 1.3), (6.2) mostra que o cálculo da coomologia, com coeficientes de torção de ordem prima com a característica, dos esquemas em grupos redutivos sobre um corpo algebricamente fechado se reduz ao problema transcendente do cálculo da coomologia dos esquemas em grupos redutivos do mesmo tipo sobre o corpo dos números complexos, cuja solução é conhecida pelos métodos da topologia algébrica, levando-se em conta o teorema de comparação (SGA 4 XI 4.4).
\end{statement}

\begin{statement}{Corolário 6.4}
Sejam $S$ um esquema e $p:G\to S$ um $S$-esquema em grupos que é extensão de um $S$-esquema em grupos finito étale $F$ por um $S$-esquema em grupos redutivo $H$. Então as conclusões de (6.2) são válidas para todo $A_S$-módulo localmente constante construtível $M$.
\end{statement}

\emph{Demonstração.} Após efetuar a mudança de base $F\to S$, pode-se supor que $F$ é da forma $\Gamma_S$, onde $\Gamma$ designa um grupo finito ordinário. Então se tem um isomorfismo de $S$-esquemas
\[
G\simeq\coprod_{\gamma\in\Gamma}H,
\]
e a conclusão segue-se de (6.2).
